\documentclass[../../main.tex]{subfiles}

\begin{document}

\chapter{Kriptografi Modern Simetris: Block Cipher}
\label{chap:block-cipher}

\begin{learningoutcome}
Setelah mempelajari bab ini, mahasiswa diharapkan mampu:
\begin{itemize}
    \item Menjelaskan konsep dasar \textit{Block Cipher} dan perbedaannya dengan \textit{Stream Cipher} (Sub-CPMK 1.2).
    \item Menganalisis berbagai mode operasi \textit{Block Cipher} (ECB, CBC, CFB, OFB, CTR) serta kelebihan dan kelemahannya (Sub-CPMK 1.2).
    \item Memahami prinsip \textit{Confusion} dan \textit{Diffusion} dalam desain algoritma kriptografi (Sub-CPMK 1.2).
    \item Menjelaskan struktur \textit{Iterated Cipher} dan mekanisme Jaringan Feistel (Sub-CPMK 1.2).
    \item Memahami proses pembangkitan kunci putaran (\textit{key scheduling}) dan pengaruh ukuran blok serta kunci terhadap keamanan (Sub-CPMK 1.2).
\end{itemize}
\end{learningoutcome}

\subfile{section-01}
\subfile{section-02}
\subfile{section-03}
\subfile{section-04}
\section{Contoh Terhitung}

\begin{example}
\textbf{ECB vs CBC.} Pakai blok mainan 4 bit: $E_K(P)$ adalah $(P\oplus K)$ yang digeser siklik ke kiri 1 bit, sedangkan $D_K(C)$ adalah hasil geser siklik ke kanan 1 bit atas $C$ yang di-XOR dengan $K$. Plainteks $1010\,0010\,0011\,1010\,1001$ (heksa $\code{A23A9}$), kunci $K=1011$ ($\code{B}$), dan $IV=\code{0000}$. Perhatikan blok ke-1 dan ke-4 sama, yaitu $A$.

\textbf{ECB} ($C_i=E_K(P_i)$, tiap blok independen):
\[
\begin{aligned}
E_K(1010)&=0010, & E_K(0010)&=0011, & E_K(0011)&=0001,\\
E_K(1010)&=0010, & E_K(1001)&=0100.
\end{aligned}
\]
sehingga $C=\code{23124}$. Blok $A$ yang berulang menghasilkan cipherteks $2$ yang sama, jadi pola plainteks tetap tampak.

\textbf{CBC} ($C_0=IV$ dan $C_i=E_K(P_i\oplus C_{i-1})$):
\[
\begin{aligned}
C_1&=E_K(1010\oplus0000)=E_K(1010)=0010,\\
C_2&=E_K(0010\oplus0010)=E_K(0000)=0111,\\
C_3&=E_K(0011\oplus0111)=E_K(0100)=1111,\\
C_4&=E_K(1010\oplus1111)=E_K(0101)=1101,\\
C_5&=E_K(1001\oplus1101)=E_K(0100)=1111,
\end{aligned}
\]
sehingga $C=\code{27FDF}$. Blok $A$ di posisi 1 dan 4 kini berubah menjadi $2$ dan $D$, jadi pengulangan tersembunyi. Dekripsi dengan $P_i=D_K(C_i)\oplus C_{i-1}$ mengembalikan plainteks semula.
\end{example}

\begin{example}
\textbf{Mode Counter (CTR).} Masih dengan kunci dan blok mainan di atas, tiap blok di-XOR dengan $E_K(\text{counter}_i)$ untuk $\text{counter}_i=i$:
\[
\begin{aligned}
E_K(0001)&=0101, & E_K(0010)&=0011, & E_K(0011)&=0001,\\
E_K(0100)&=1111, & E_K(0101)&=1101.
\end{aligned}
\]
Kemudian $C_i=P_i\oplus E_K(\text{counter}_i)$:
\[
C_1=1010\oplus0101=1111,\quad C_2=0010\oplus0011=0001,\quad C_3=0011\oplus0001=0010,
\]
\[
C_4=1010\oplus1111=0101,\quad C_5=1001\oplus1101=0100,\qquad C=\code{F1254}.
\]
Karena setiap blok hanya bergantung pada counternya, bukan pada cipherteks sebelumnya, blok dapat dienkripsi paralel dan diakses acak, berbeda dari CBC.
\end{example}


\begin{obeactivity}
\textbf{Aktivitas 6.1: Analisis Mode ECB vs CBC}
1. Gunakan alat bantu enkripsi (seperti CyberChef) untuk mengenkripsi sebuah gambar sederhana (misal: gambar logo kecil) menggunakan AES mode ECB dan AES mode CBC.
2. Amati hasil cipherteks yang dihasilkan. Apakah gambar asli masih terlihat pada mode ECB?
3. Diskusikan mengapa mode CBC lebih aman untuk data yang memiliki pola berulang.
\end{obeactivity}

Langkah-langkah pada Aktivitas 6.1 dapat dijalankan langsung melalui program
pada Listing~\ref{lst:mode-operasi}. Program itu memakai block cipher mainan
berukuran 4 byte untuk membandingkan ECB dan CBC pada plainteks berpola
berulang, sehingga kebocoran pola pada ECB terlihat jelas tanpa perlu mengolah
gambar.

\lstinputlisting[
    language=Python,
    caption={Mode operasi ECB dan CBC: kebocoran pola pada blok yang berulang},
    label={lst:mode-operasi}
]{\codefile{python/mode_operasi.py}}

Program yang sama tersedia dalam C++ (\texttt{code/cpp/mode\_operasi.cpp}) dan
MATLAB (\texttt{code/matlab/mode\_operasi.m}).


\begin{obereflection}
\textbf{Latihan 6.1}
1. Jika sebuah blok cipher memiliki ukuran 64 bit dan kita menggunakan mode ECB, apa yang terjadi jika dua blok plainteks identik dikirimkan?
2. Jelaskan perbedaan mendasar antara peran S-Box dan P-Box dalam mencapai \textit{confusion} dan \textit{diffusion}.
3. Mengapa Jaringan Feistel memudahkan implementasi dekripsi?
\end{obereflection}

\section{Rangkuman}
\begin{itemize}
    \item \textit{Block cipher} membagi plainteks menjadi blok tetap $n$ bit ($C=E_K(P)$, $P=D_K(C)$); panjang cipherteks sama dengan panjang plainteks, sedangkan panjang kunci eksternal bebas.
    \item Untuk setiap kunci, $E_K$ harus merupakan permutasi bijektif atas himpunan $2^n$ blok; cipher blok hanya memilih $2^m$ permutasi dari $2^n!$ permutasi yang tersedia.
    \item Bila pesan tidak habis dibagi ukuran blok, blok terakhir ditambahi bit \textit{padding} agar panjangnya kembali $n$ bit.
    \item Confusion menyembunyikan hubungan kunci--cipherteks melalui substitusi non-linear (S-box, misalnya $S_1$ DES: $\code{110100}\to\code{1001}$), sedangkan diffusion menyebarkan pengaruh satu bit plainteks ke banyak bit cipherteks melalui permutasi.
    \item Confusion dan diffusion baru tercapai menyeluruh setelah substitusi dan permutasi diulang berkali-kali sebagai \textit{iterated cipher}, dengan kunci putaran yang berbeda pada setiap putaran.
    \item Mode ECB ($C_i=E_K(P_i)$) sederhana dan mendukung akses acak, tetapi blok plainteks yang sama selalu menghasilkan cipherteks yang sama sehingga pola data bocor dan blok dapat dibuang atau disusun ulang.
    \item Mode CBC ($C_i=E_K(P_i\oplus C_{i-1})$, $C_0=IV$) menghapus pola berulang dan menutup manipulasi blok, namun galat satu bit menjalar ke seluruh blok cipherteks berikutnya saat enkripsi.
    \item Mode CFB ($C_j=P_j\oplus\mathrm{MSB}_s(E_K(X_j))$) memperlakukan cipher blok seperti cipher alir untuk data yang belum genap satu blok, dengan umpan balik diambil dari cipherteks.
    \item Mode OFB mengambil umpan balik dari keluaran $E_K$, bukan dari cipherteks, sehingga aliran kuncinya tidak bergantung pada data dan galat tidak menjalar sama sekali.
    \item Mode CTR ($C_i=P_i\oplus E_K(\text{counter}_i)$) mendukung pemrosesan paralel penuh dan akses acak karena tiap blok tidak bergantung pada blok lain; nilai counter wajib tidak pernah berulang.
    \item Jaringan Feistel ($L_i=R_{i-1}$, $R_i=L_{i-1}\oplus f(R_{i-1},K_i)$) bersifat reversibel sehingga dekripsi memakai urutan kunci terbalik tanpa fungsi $f^{-1}$ terpisah; sifat ini tidak bergantung pada bentuk fungsi $f$.
    \item Jaringan SPN seperti AES menuntut setiap komponen putaran dapat dibalik, sehingga inversi S-box dan matriks pencampur harus dirancang khusus.
    \item \textit{Key scheduling} menurunkan kunci eksternal menjadi $r$ kunci putaran yang berbeda; DES mengubah kunci 64 bit menjadi 16 kunci putaran 48 bit.
    \item Ukuran blok menentukan luas diffusion, sedangkan ukuran kunci menentukan ketahanan terhadap pencarian menyeluruh; blok 64 bit perlu diganti kunci sebelum $2^{32}$ blok karena paradoks ulang tahun.
\end{itemize}


\begin{obeassessment}
\textbf{Kuis Bab 6}
1. Berikan contoh skenario di mana mode CTR lebih menguntungkan dibandingkan mode CBC.
2. Apa yang dimaksud dengan \textit{error propagation} pada mode CBC dan bagaimana mode OFB mengatasinya?
3. Mengapa penggunaan IV (\textit{Initialization Vector}) sangat penting pada mode CBC?
4. Jelaskan mengapa mode ECB dapat dimanipulasi oleh pihak lawan untuk mengubah makna pesan meskipun kuncinya tidak diketahui.
5. Sebuah cipher blok berukuran 4 bit memakai $E_K(P)=(P\oplus K)\ll 1$ dengan $K=\code{1011}$. Hitung cipherteks dari plainteks $\code{0110}$ pada mode ECB.
6. Pada mode OFB, mengapa aliran kunci dapat dihitung lebih dahulu sebelum plainteks diketahui? Sebutkan satu keuntungan praktis dari sifat itu.
7. Jelaskan hubungan antara jumlah putaran sebuah \textit{iterated cipher} dengan tingkat keamanan dan dengan biaya komputasi. Mengapa jumlah putaran tidak dibuat sebesar mungkin?
8. Apa perbedaan utama antara jaringan Feistel dan jaringan SPN dari sudut pandang perancang algoritma?
\end{obeassessment}
\begin{competencychecklist}
\begin{tabularx}{\linewidth}{|X|c|}
\hline
Indikator Kompetensi & Tercapai \\
\hline
Saya dapat menjelaskan perbedaan \textit{block cipher} dan \textit{stream cipher} & [ ] \\
\hline
Saya dapat menjelaskan mengapa $E_K$ harus merupakan permutasi bijektif & [ ] \\
\hline
Saya dapat menangani blok terakhir yang tidak lengkap dengan \textit{padding} & [ ] \\
\hline
Saya dapat membedakan mode operasi ECB, CBC, CFB, OFB, dan CTR & [ ] \\
\hline
Saya dapat menuliskan persamaan enkripsi dan dekripsi tiap mode operasi & [ ] \\
\hline
Saya dapat menganalisis penjalaran galat pada mode CBC, CFB, dan OFB & [ ] \\
\hline
Saya dapat menjelaskan peran IV pada CBC dan peran counter pada CTR & [ ] \\
\hline
Saya dapat menjelaskan prinsip \textit{confusion} dan \textit{diffusion} & [ ] \\
\hline
Saya dapat menghitung substitusi satu masukan pada S-box DES & [ ] \\
\hline
Saya dapat menjelaskan efek \textit{avalanche} pada cipher blok & [ ] \\
\hline
Saya dapat menguraikan mekanisme kerja Jaringan Feistel & [ ] \\
\hline
Saya dapat membuktikan sifat reversibel Jaringan Feistel & [ ] \\
\hline
Saya dapat menjelaskan tujuan \textit{key scheduling} & [ ] \\
\hline
Saya dapat menganalisis pengaruh panjang kunci terhadap keamanan sistem & [ ] \\
\hline
\end{tabularx}
\end{competencychecklist}


\end{document}
